\section{Introduction}
Large language models (LLMs) have evolved from basic systems into autonomous agents capable of complex planning, reasoning, and tool execution \cite{agent_survey, autogpt_p}. By breaking down problems and interacting with external APIs, these agents offer immense practical value. However, as these systems advance, they face a major challenge: severe operational costs. The constant cycle of reasoning and action consumes massive context window space, driving up token usage and latency. Therefore, a primary focus of current research is building frameworks that maintain high performance while drastically reducing computational costs.\\

\section{Cost and Memory Limitations in Cloud Agents}
Many popular tools, such as LangChain, LangGraph, and LangSmith, help manage how AI agents execute tasks step by step \cite{langchain_langgraph}. However, cloud-based systems today, including Claude Code \cite{claude_code_dive}, still rely on massive context windows and generate temporary code for every single task, making the system slower and greatly increasing operational costs. To address high costs, studies suggest treating AI tokens as a limited and valuable resource \cite{token_economics, ai_tokenomics, marginal_token}. By carefully managing tokens, developers can build fast, scalable agents while keeping costs low \cite{efficient_agents}. A critical component of cost reduction is long-term memory. Studies comparing databases like Chroma, Qdrant, and FAISS \cite{chroma_faiss_analysis} show that choosing the right local vector storage helps the agent find information quickly, preventing it from repeating expensive cloud requests. Additionally, lightweight tools like FLAML \cite{flaml} can be integrated into these systems to save computational money \cite{cost_aware_automl}.\\

\section{Challenges in Traditional Agent Architectures}
A foundational design for many modern AI agents is the ReAct (Reasoning + Acting) model \cite{yao2023react}. In this model, the agent thinks about a problem, takes an action, and then checks the result. While reliable, this back-and-forth process has a major flaw: as the agent takes more steps, the history of its actions grows very long. This causes the system to use more tokens, become slower, and cost more money \cite{token_economics, efficient_agents}. Furthermore, standard agents usually handle tool calling and memory poorly \cite{efficient_tool_select, langchain_langgraph}. Most agents just append the entire conversation history to their active memory, relying on massive logs that waste tokens rather than properly saving project states \cite{ai_agent-memory}. Another major challenge is how current agents search for information. Instead of using a unified semantic memory system, many agents rely on disjointed storage formats, such as raw Markdown files, plain text files, or basic SQL databases. To find the root cause of a problem, these systems frequently use brute-force regex searches or generic file-finding commands. Each time they perform a search, they must send massive blocks of unstructured text back to the cloud LLM to interpret the results. This creates an endless loop of searching, reading, and re-prompting the LLM, which wastes significant time and operational budget \cite{claude_code_dive, ai_agent-memory}. These inefficiencies are highly visible in traditional specialized agents. To solve difficult problems, systems use specific roles like Code Agents, Browser Agents, OS Agents, and Testing Agents. However, Code agents like SWE-agent \cite{yang2024sweagent} often forget file locations and repeatedly search the system. Browser agents like BrowserGym \cite{chezelles2025browsergym} constantly re-read web pages instead of remembering them. OS agents like OSWorld \cite{xie2024osworld} easily lose track of their environment between steps. Finally, testing agents often dump huge error logs directly into the AI, unnecessarily wasting token space and increasing costs.

\section{Research Gap}
Despite recent improvements, today's popular AI agents still rely too much on expensive cloud services and struggle to be efficient \cite{efficient_agents}. Specifically, there are five major gaps in current systems:

\begin{itemize}
    \item \textbf{No Long-Term Memory for Workflows:} Current agents easily forget how they solved a problem \cite{ai_agent-memory}. They lack a method to instantly remember and reuse tools for finished tasks.
    \item \textbf{Repeating the Same Work:} Because agents do not save step-by-step action plans, they waste money re-solving the exact same problems from scratch whenever a user asks a similar question \cite{claude_code_dive, token_economics}.
    \item \textbf{High Costs from Overusing Cloud AI:} Existing systems send almost every request directly to massive cloud models \cite{ai_tokenomics}. There is a lack of systems attempting to solve simple problems locally first.
    \item \textbf{Missing a Step-by-Step Filter (Triage):} Current agents lack a strict filtering process \cite{marginal_token}. There is no system that tries fast local processing first, then a tiny local AI, and only pays for an expensive cloud AI as a last resort.
    \item \textbf{Impractical Fine-Tuning:} Many systems attempt to solve memory issues by fine-tuning local models \cite{agent_survey}. However, fine-tuning is expensive, static, and cannot learn new tool APIs dynamically. Experiential learning via local databases is required instead of retraining.
\end{itemize}

This research aims to fix these problems by building the Ecogent architecture. By focusing on local processing, smart memory, and reusing past work, Ecogent provides an affordable and sustainable solution for the IT industry.
\begin{table*}[htbp]
\centering
\caption{Comparison of Recent Cost- and Token-Efficient Agent Research}
\label{tab:related_work_cost_efficient_agents}
\begin{tabular}{c p{0.28\textwidth} p{0.28\textwidth} p{0.28\textwidth}}
\hline
\textbf{Ref.} & \textbf{Methodology} & \textbf{Result} & \textbf{Scope for Improvement} \\
\hline

\cite{token_economics} &
Survey-based ``Token Economics'' framework modeling tokens as resources across single-agent, multi-agent, and ecosystem levels. Studies routing, memory, and budget allocation.
&
Identifies token consumption and communication as major costs. Recommends routing, caching, and budget-aware reasoning to improve cost-quality efficiency.
&
Lacks a unified benchmark. Differing cost definitions across the reviewed studies make direct comparison difficult.
\\
\hline

\cite{ai_tokenomics}&
Introduces AgentAssay for regression testing using SPRT, behavioral fingerprinting, and trace-first analysis. Evaluated over 7,605 trials.
&
Reports 5.3$\times$ variation in token generation. Trace-first testing achieved 100\% cost savings when existing traces were available.
&
Only 5 models and 3 scenarios were tested. Broad multi-agent settings and physical-world tasks were not evaluated.
\\
\hline

\cite{marginal_token} &
Proposes marginal token allocation across model routing, agent actions, inference serving, and post-training, balancing token cost, latency, and risk.
&
Theoretically shows that simply minimizing tokens causes system congestion and cache misuse. Proposes a shared price signal.
&
No empirical validation. The framework assumes latency and risk can be combined into comparable costs, which may fail in complex tasks.
\\
\hline

\cite{efficient_agents} &
Empirically studies LLM backbone, planning, tools, and memory on the GAIA benchmark to build a task-efficient configuration.
&
Retained 96.7\% of OWL's performance while reducing cost from \$0.398 to \$0.228 and improving cost-of-pass by 28.4\%.
&
Evaluation is mainly based on GAIA. Results may not generalize to long-horizon, domain-specific, or multi-agent workloads.
\\
\hline

\cite{cost_aware_automl} &
Presents a cost-aware AutoML agent using generative reasoning, cost evaluation, and dynamic search pruning, balanced by a utility function.
&
Achieved 25--60\% cost reduction on 75 datasets while maintaining or improving performance over Auto-sklearn and H2O AutoML.
&
Focused primarily on AutoML tasks. High dependency on the specific utility function makes generalization to other LLM agents uncertain.
\\
\hline

\end{tabular}
\end{table*}